\beginsong{La Bohème}[by={Charles Aznavour}]
\beginverse
\[Fm]Je vous parle d'un temps que les moins de vingt ans
Ne peuvent pas con\[Cm]naître. Montmartre en ce temps-\[Fm]là
Accrochait ses lilas jusque sous nos fe\[Cm]nêtres.
Et si l'humble gar\[Fm]ni qui nous servait de nid
Ne payait pas de \[Cm]mine, c'est là qu'on s'est con\[Fm]nus,
Moi qui criait fa\[G9]mine et toi qui posais \[Cm]nue.
\endverse
\beginchorus
\[C7]La bo\[Fm]hème, la bo\[Cm]hème,
Ça vou\[Fm]lait dire on est heu\[G7]reux. \[Cm] \[C7]
La bo\[Fm]hème, la bo\[Cm]hème,
Nous ne man\[Fm]gions qu'un jour sur \[G7]deux. \[Cm]
\endchorus
\beginverse
\[Fm]Dans les cafés voisins nous étions quelques-uns
Qui attendions la \[Cm]gloire, et bien que misé\[Fm]reux
Avec le ventre creux nous ne cessions d'y \[Cm]croire.
Et quand quelque bis\[Fm]tro contre un bon repas chaud
Nous prenait une \[Cm]toile, nous récitions des \[Fm]vers
Groupés autour du \[G9]poêle en oubliant l'hi\[Cm]ver.
\endverse
\beginverse
\[Fm]Souvent il m'arrivait devant mon chevalet
De passer des nuits \[Cm]blanches, retouchant le des\[Fm]sin
De la ligne d'un sein, du galbe d'une \[Cm]hanche.
Et ce n'est qu'au ma\[Fm]tin qu'on s'asseyait enfin
Devant un café-\[Cm]crème, épuisés mais ra\[Fm]vis,
Fallait-il que l'on \[G9]s'aime et qu'on aime la \[Cm]vie.
\endverse
\beginverse
\[Fm]Quand au hasard des jours je m'en vais faire un tour
À mon ancienne a\[Cm]dresse, je ne reconnais \[Fm]plus
Ni les murs ni les rues qui ont vu ma jeu\[Cm]nesse.
En haut d'un esca\[Fm]lier je cherche l'atelier
Dont plus rien ne sub\[Cm]siste, dans son nouveau dé\[Fm]cor
Montmartre semble \[G9]triste et les lilas sont \[Cm]morts.
\endverse
\endsong
